% Physical unit 029; accepted Session 03 after fresh 400-dpi authority replay.
\GordanUnitStart{029}{printed page 24}
\GordanTopFolio{24}
\begingroup\small
\[
\begin{gathered}
 \tg T={\cot l_0\over Su},\qquad
 \sin T={1\over \sqrt{1+\tg^2l_0\,S^2u}},\qquad
 \cos T={Su\over \sqrt{\cot^2l_0+S^2u}},\\
 s={b\over x'}\int_0^u \Delta^2u\,du,\\
 \psi={\cot l_0\over \sin B_0}\int_0^u du\,{\Delta^2u\over \cot^2l_0+S^2u}.
\end{gathered}
\]

\sect{3}

In this way the trigonometric functions of the angles
\[
 B\qquad l\qquad T
\]
are represented by elliptic functions of the quantity $u$; it remains to determine, by means of its $\Theta$ functions, the extension and the longitude.

For this reason the formula for the quantity $s$ is transformed into the following:
\[
 {x's\over b}=u-\int_0^u x^2S^2u\,du
\]
in which, after using the formula proposed in the \emph{Fundamenta}, p. 133, and putting
\[
 \int_0^{\frac{\pi}{2}}{d\varphi\over \sqrt{1-x^2\sin^2\varphi}}=K,
 \qquad
 \int_0^{\frac{\pi}{2}}d\varphi\sqrt{1-x^2\sin^2\varphi}=E,
\]
one obtains:
\[
 {2Kx'\over b\pi}\cdot s={2K\over\pi}u+{2K\over\pi}Zu-u\left({2K\over\pi}-{2E\over\pi}\right)
 ={2E\over\pi}u+{2K\over\pi}Zu
\]
\[
 s={b\over x'}\cdot {E\over K}u+{b\over x'}Zu.
\]
And if, in place of the constant $E$, the constant $\Sigma(l_0)$ is introduced, into which $s$ passes on putting $u=K$; that is, if we put
\[
 {b\over x'}E. = \Sigma(l_0),
\]
then
\[
 s-{u\over K}\Sigma(l_0)={b\over x'}Zu.
\]

\sect{4}

It is more difficult to refer the longitude to $\Theta$ functions. For this purpose the method used by the illustrious Jacobi in his booklet
\endgroup
